% ============================================================================
% Section II: Literature Review and State of the Art
% ============================================================================

\section{Literature Review}
\label{sec:literature}

\subsection{Generation and Transmission Expansion Planning}

GTEP is a classical problem in power systems that has been studied
extensively since the 1970s.  Romero et al.~\cite{romero2002test} provided
canonical test systems and mathematical models---including transportation,
hybrid, DC power flow, and disjunctive formulations---that remain standard
benchmarks.  Hemmati et al.~\cite{hemmati2013comprehensive} presented a
comprehensive review covering generation expansion planning (GEP),
transmission expansion planning (TEP), and their integration.  More
recently, Mahdavi et al.~\cite{mahdavi2019tep_survey} classified TEP
literature into categories by mathematical model, solution method, and
uncertainty handling, identifying a trend toward stochastic,
multi-objective, and market-based formulations.

Lumbreras and Ramos~\cite{lumbreras2016progressive} surveyed modern
challenges in TEP including renewable integration, energy storage, and
the need for multi-stage stochastic models.  These challenges motivate
the development of tools that can handle joint generation, transmission,
and storage expansion under operational uncertainty.

\subsection{DC Optimal Power Flow}

The DC power flow approximation linearizes the AC power flow equations
by assuming lossless lines, flat voltage profiles, and small angle
differences~\cite{stott2009dc_opf}.  Despite its simplifications, DC OPF
is the standard formulation for TEP because it yields a linear program
amenable to efficient solution~\cite{frank2016or_opf}.  Molzahn and
Hiskens~\cite{molzahn2019opf_survey} provided a comprehensive survey of
relaxations and approximations, confirming that DC OPF remains the
workhorse for expansion planning despite advances in convex relaxations
of AC OPF.

\gtopt{} adopts the $B$-$\theta$ formulation (susceptance-angle) for
DC OPF, as recommended by Stott et al.~\cite{stott2009dc_opf}, with
extensions for resistive losses via piecewise-linear approximations and
phase-shifting transformers.

\subsection{Stochastic Dual Dynamic Programming}

The Stochastic Dual Dynamic Programming (SDDP) algorithm, introduced by
Pereira and Pinto~\cite{pereira1991sddp}, is the standard method for
solving large-scale multi-stage stochastic programs in hydrothermal
scheduling.  The algorithm constructs piecewise-linear approximations of
the expected cost-to-go function through iterative forward and backward
passes.

Shapiro~\cite{shapiro2011sddp_analysis} provided theoretical convergence
analysis of SDDP, and Philpott and de~Matos~\cite{philpott2013sddp_convergence}
extended the algorithm with risk-averse sampling.  Recent advances include
batch learning SDDP~\cite{leclercq2023batch_sddp}, which frames the
algorithm as Q-learning for improved convergence, and extensions to
stochastic OPF under uncertainty~\cite{fletcher2024sddp_opf}.

In the Julia ecosystem, Dowson and Kapelevich~\cite{dowson2020sddp_jl}
developed SDDP.jl, a modular SDDP implementation.  \gtopt{} provides a
C++ SDDP implementation with support for multi-scenario apertures, elastic
filters for feasibility management, and boundary cost cuts for rolling
horizon studies.

\subsection{Energy Storage in Expansion Planning}

The integration of battery energy storage systems (BESS) into expansion
planning has received significant attention.  Go et al.~\cite{Go2016ess_expansion}
assessed the economic value of co-optimized grid-scale storage investments.
Frazier et al.~\cite{frazier2023bess_milp} presented an integrated MILP
for generation, transmission, and storage expansion with Benders
decomposition.

Sandia National Laboratories released QuESt Planning~\cite{sandia2024quest},
an open-source capacity expansion tool focused on storage, while
IDAES-GTEP~\cite{idaes_gtep2024} provides a modular Pyomo-based framework
for GTEP.  \gtopt{} supports BESS with detailed charge/discharge efficiency
modeling, state-of-charge tracking across stages, and co-located
generation-storage coupling via its converter mechanism.

\subsection{Open-Source Power System Tools}

Table~\ref{tab:tool_comparison} compares \gtopt{} with the leading
open-source GTEP tools.  PyPSA~\cite{brown2018pypsa} is the most widely
adopted, offering multi-sector coupling and a large community, but uses
Python for LP assembly which can bottleneck on large problems.
GenX~\cite{jenkins2017genx} provides a flexible Julia/JuMP framework but
lacks native DC OPF and hydro cascade modeling.  pandapower~\cite{thurner2018pandapower}
excels at power flow analysis but does not support expansion planning.

\begin{table*}[t]
  \centering
  \caption{Comparison of Open-Source GTEP and Power System Tools}
  \label{tab:tool_comparison}
  \begin{tabular}{@{}lcccccc@{}}
    \toprule
    \textbf{Feature} & \textbf{gtopt} & \textbf{PyPSA} & \textbf{GenX}
      & \textbf{IDAES-GTEP} & \textbf{QuESt} & \textbf{PLP} \\
    \midrule
    Language          & C++26  & Python  & Julia   & Python  & Python  & Fortran \\
    License           & BSD-3  & MIT     & GPL-2   & BSD-3   & BSD-3   & Proprietary \\
    DC OPF (Kirchhoff)& \checkmark & \checkmark & Transport & \checkmark & -- & \checkmark \\
    Hydro cascades    & \checkmark & Partial & -- & -- & -- & \checkmark \\
    Battery storage   & \checkmark & \checkmark & \checkmark & Partial & \checkmark & Partial \\
    SDDP solver       & \checkmark & -- & -- & -- & -- & \checkmark \\
    Capacity expansion& \checkmark & \checkmark & \checkmark & \checkmark & \checkmark & -- \\
    Multi-scenario    & \checkmark & Partial & Partial & -- & -- & \checkmark \\
    Multi-sector      & -- & \checkmark & -- & -- & -- & -- \\
    Parquet I/O       & \checkmark & -- & -- & -- & -- & -- \\
    Web GUI           & \checkmark & -- & -- & -- & \checkmark & -- \\
    \bottomrule
  \end{tabular}
\end{table*}

\gtopt{} distinguishes itself through: (i)~C++ sparse-matrix assembly
for high performance, (ii)~complete hydro cascade modeling with
reservoirs, waterways, turbines, and filtration, (iii)~three solution
methods including SDDP and cascade decomposition, (iv)~native Apache
Parquet I/O for large time-series, and (v)~a comprehensive Python
toolchain for data preparation and cross-validation.
